% Proposed replacement for the unresolved abelian-type portion of C.6.
% Prepared by Part D; Part C owns integration into sections/part_c.tex.
\begin{lemma}[A generator with potentially good reduction]
\label{c:potentially-good-generator}
Let $M$ be a rational motive of abelian type over a number field $K$, in either
the absolute-Hodge or motivated category.  Let $v$ have residue characteristic
$p$.  If $M_{\ell_0}$ is potentially unramified at $v$ for one $\ell_0\ne p$,
then, after a finite Galois extension $E/K$, the object $M_E$ belongs to the
tensor category generated by finitely many abelian varieties over $E$ which
have good reduction at every place above $v$.
\end{lemma}

\begin{proof}
We first work over $\overline\Q$, using the chosen complex embedding.
The weak Mumford--Tate lifting construction of
\citet[Theorem 2.12 and Corollary 2.15, pp.~15--20]{c_noot2006} gives an
abelian variety $B$ with a central isogeny
\[
 \pi^{\mathrm{der}}:\operatorname{MT}(B)^{\mathrm{der}}
       \longrightarrow\operatorname{MT}(M)^{\mathrm{der}}
\]
identifying their adjoint Hodge homomorphisms.  Here descent to
$\overline\Q$ is part of the argument, not an assumption: start with an
abelian generator $A/\overline\Q$ of $M$, take its Mumford--Tate decomposed
weak lift $B_0$, and descend $B_0$ by Noot's Proposition 3.1 and Corollary 3.2,
pp.~20--21.  The proof of Corollary 2.15 selects the factors of $B_0$ whose
$\Q$-simple derived groups survive in $\operatorname{MT}(M)$.  These factors
descend up to isogeny because their rational endomorphism idempotents descend
under an algebraically closed field extension.  Rational motives only require
this isogeny-level descent.

The proof of Noot's Proposition 3.4, pp.~22--23, applies with the Hodge
realization of $M$ in place of $H^1(A)$ and gives
\begin{equation}
\label{c:weak-lift-generates}
 M_{\overline\Q}\ \in\
 \langle h^1(B),\text{CM motives}\rangle^{\otimes}.
\end{equation}
To make this extension of the source statement explicit, use the common
group formed from the central tori and $\operatorname{MT}(B)^{\mathrm{der}}$.
Its representation on $\omega_B(M)$ is generated by the faithful
representation on $H^1(B)$ and representations of its torus quotient.
For an irreducible constituent $W_2$ of $\omega_B(M)$, choose an irreducible
constituent $W_1$ of a tensor construction on $H^1(B)$ and a torus
representation $W_3$ with $W_2\subset W_3\otimes W_1$.
The nonzero contraction $W_2\otimes W_1^\vee\to W_3$ shows that $W_3$ itself
is an actual polarizable Hodge subquotient, of CM type, with integral
Hodge degrees.  Full faithfulness of Betti realization for abelian motives
and the realization of CM Hodge structures by CM abelian varieties prove
\eqref{c:weak-lift-generates}.  In the motivated category the same step uses
Andr\'e's Theorem 0.6.2.  Only finitely many CM varieties and morphisms are
needed for this particular $M$, so all these data descend to a finite
extension of $K$.

Put $H=\operatorname{MT}(M_{\overline\Q}\oplus h^1(B))$.
Its two projections have the same adjoint Hodge datum.  The Hodge
homomorphism defining $H$ lies in the fibre product over that adjoint group,
and projects onto both Mumford--Tate groups.  The kernel of the projection
$H\to\operatorname{MT}(M)$ lies in $\{1\}\times Z(\operatorname{MT}(B))$,
so is central in $H$ and meets its semisimple derived group in a finite
group.  Thus $H^{\mathrm{der}}\to\operatorname{MT}(M)^{\mathrm{der}}$ is a
central isogeny.  After a finite extension, the arithmetic realizations of
$M\oplus h^1(B)$ take values in $H(\Q_\ell)$ for every $\ell$: finitely many
absolute Hodge tensors cut out $H$, and one extension defines all of them.

Fix a place above $v$ and take a further finite local extension over which
$M_{\ell_0}$ is unramified and $B$ is semistable.  An inertia element acts on
$M_{\ell_0}\oplus H^1(B,\Q_{\ell_0})$ as
\[
 u=(1,\rho_{B,\ell_0}(\tau))\in H(\Q_{\ell_0}),
\]
and $u$ is unipotent by the semistable monodromy theorem.  Its image in the
torus $H/H^{\mathrm{der}}$ is trivial, so $u\in H^{\mathrm{der}}$.
Its image in $\operatorname{MT}(M)^{\mathrm{der}}$ is trivial, and hence
$u$ lies in a finite kernel.  A finite-order unipotent in characteristic
zero is the identity.  Thus $H^1(B,\Q_{\ell_0})$ is unramified, and
N\'eron--Ogg--Shafarevich gives good reduction of $B$ over this local
extension.  The same reasoning applies at \emph{every} place above $v$,
since potential unramifiedness of $M_{\ell_0}$ transfers to all those
completions.

The auxiliary CM abelian varieties are potentially good as well: after
fixing their Hodge tensors and taking a semistable extension, their inertia
is both unipotent and contained in a torus, hence trivial.  The same
criterion applies.  There are only finitely many varieties and places.
Choose finite extensions of the completions giving good reduction and
realize them within one finite global extension by weak approximation and
Krasner's lemma.  Pass to its normal closure over $K$; good reduction is
preserved by further extension.  This gives the asserted $E$ and proves
the direction of transfer required in the lemma.
\end{proof}

\begin{theorem}[Independence of unramifiedness for all abelian-type motives]
\label{c:ell-independence-abelian}
\status{Proved, using the October 2026 rationality theorem}
Let $M$ be a rational absolute-Hodge or motivated motive of abelian type
over a number field $K$, and let $v$ have residue characteristic $p$.
If $M_{\ell_0}$ is unramified at $v$ for one $\ell_0\ne p$, then
$M_\ell$ is unramified at $v$ for every $\ell\ne p$.
The statement includes $p=2$, and includes $\ell=2$ when $p\ne2$.
\end{theorem}

\begin{proof}
Apply Lemma~\ref{c:potentially-good-generator}.  Over a finite Galois
extension $E/K$, the motive $M_E$ is a summand of a finite direct sum $P$
of motives $h^{i_a}(A_a)(m_a)$, where the $A_a$ are powers or products of
abelian varieties with good reduction at \emph{every} place above $v$.
Indeed, tensor constructions on $h^1$, their duals and Tate twists occur
as K\"unneth summands in such powers.  A finite direct sum permits different
degrees $i_a$ and twists $m_a$.

Restriction and induction give morphisms
\[
 M\longrightarrow\operatorname{Ind}_E^K(M_E)\longrightarrow M
\]
whose composite is $[E:K]$.  After division by this nonzero rational
integer, $M$ is a summand of $\operatorname{Ind}_E^K(P)$.
These are finite \'etale correspondence maps in either motivic category;
after extension to $E$ they are the diagonal and summation maps over
$\Gal(E/K)$.  Thus the descent assertion is established over the original
$K$, without inferring it from trivial inertia over a ramified extension.

Write $e$ for the resulting rational Hodge, equivalently motivated,
projector.  The induced ambient varieties are finite disjoint unions of
the conjugate $A_a$.  Fix the chosen place of $\overline\Q$ above $v$ and
take a common finite Galois local extension $L/K_v$ splitting the finite
extensions of induction and giving good models for all components.
Smooth proper base change shows, simultaneously for every $\ell\ne p$,
that inertia on the ambient object factors through the same finite
quotient
\[
 \Delta=I_v/(I_v\cap G_L).
\]
It follows that this also holds on $M_\ell$.  Tate twists are trivial on
this inertia group when $\ell\ne p$.

By uniqueness of the abelian scheme models, the descent maps extend to
them.  On the geometric special fibres, each $\tau\in\Delta$ acts by a
permutation of components and origin-preserving algebraic isomorphisms
$u:X_0\to Y_0$ between conjugate abelian varieties.  Inertia fixes the
algebraic residue field, so these maps are actual homomorphisms over that
field.  Their pullbacks describe the cohomological inertia action.

The graph of every such $u$ is a rational Lefschetz class on $X_0\times Y_0$.
If $d=\dim Y_0=\dim X_0$ and $f(x,y)=y-u(x)$, then in every Weil
cohomology,
\begin{equation}
\label{c:inertia-graph-lefschetz}
 [\Gamma_u]=f^*[0_{Y_0}]
   =f^*\left(\frac{\theta_{Y_0}^{d}}
                         {\deg(\theta_{Y_0}^{d})}\right),
\end{equation}
where $\theta_{Y_0}$ is any ample divisor class.  This uses the common
degree normalization of top classes.  The expression on the right is one
rational polynomial in divisor classes, with the same cycle realizations
for all $\ell$.

Use the correspondence convention
$Z_*(x)=(\mathrm{pr}_2)_*([Z]\smile\mathrm{pr}_1^*x)$.
Then $\Gamma_u$ acts by $u_*=(u^{-1})^*$, and the alternating trace of
$e_*\circ u^*$ is the intersection pairing of $e$ with $\Gamma_u$.
For a projector supported in degree $i$, the ordinary trace is therefore
$(-1)^i$ times that pairing.  On a disjoint union, the same identity pairs
the matrix entries of $e$ with the graphs indexed by the component
permutation.  For a direct sum of different degrees and twists, only
the diagonal blocks relative to those Galois-stable summands contribute
to the trace.  Consequently
\[
 \operatorname{Tr}(\tau\mid M_\ell)
\]
is a finite sum, with the indicated signs, of pairings of specialized
rational Hodge classes with complementary rational Lefschetz classes on
products of abelian varieties with good reduction.  The generic matrix
entries of $e$ are Hodge classes on these products; realization and smooth
proper specialization preserve their actions, tensor operations, duality
and K\"unneth decompositions.

The rationality theorem \citep[Theorem 1.1]{c_rationality2026}, also
\citep[Theorem 1.1, pp.~1--2]{c_rationalityexposition2026}, says that every
such pairing is one rational number in every $\ell$-adic realization
$\ell\ne p$, including residue characteristic $2$.  Hence for each
$\tau\in\Delta$ there is a common rational trace $t_\tau$.
The assumed unramifiedness at $\ell_0$ gives
$t_\tau=\dim_\Q\omega_B(M)$.  In every other realization the operator
$\tau$ has finite order.  Its eigenvalues are roots of unity, and under a
complex embedding their sum equals their number.  Equality in the triangle
inequality forces every eigenvalue to be $1$.  Finite-order operators in
characteristic zero are semisimple, so $\tau$ is the identity on $M_\ell$.
This holds for every $\tau$, proving the theorem.
\end{proof}

\begin{remark}[The exact role of the new input]
The proof does not apply N\'eron--Ogg--Shafarevich to an arbitrary original
generator: such a generator need not have good reduction when its summand
is unramified.  Lemma~\ref{c:potentially-good-generator} changes the
generator and proves its potential good reduction.  The finite descent
action is then recovered over $K$ by the common trace pairing.  Since
\eqref{c:inertia-graph-lefschetz} is Lefschetz, the proof uses the
rationality theorem alone; it does not use the CM Hodge theorem or assume
algebraicity of the Hodge projector.  It proves independence of
unramifiedness and, for potentially unramified objects, independence of
the finite inertia character.  It does not assert full Weil--Deligne
compatibility for objects with nonzero monodromy.
\end{remark}
